\section{司法制度和法律职业道德概述}
\begin{enumerate}
  \item 8 - 10 分。一部分常识判断，两诉知识，专门学习，运气。
  \item 司法机关。题目没有说明主体，偏向法院。
  \item 选项出现名人名言，一律对。
  \item \zd{司法特征}。\zd{司法机关运用法律}。vs 行政活动 vs 社会活动。行政偏直接，司法偏间接。
        \begin{enumerate}
          \item （行使职权）独立性。系统、个体独立。检察系统是整体独立，不包括个体独立。
          \item 被动性。不告（请求）不理。
          \item 交涉性。多方参与。
          \item 程序性。
          \item 普遍性。最普适的纠纷处理途径。
          \item 终极/局性。公信力。
        \end{enumerate}
  \item 功能。
        \begin{enumerate}
          \item 应然/实然。应该/实际怎样。理想/现实。
          \item \zd{直接 \yjt 解决纠纷}。间接功能，不是直接的就是间接的，大概率 ABCD。
        \end{enumerate}
  \item 司法制度。审判、检察。广义：一切跟司法有关的。我国：
        \begin{enumerate}
          \item 4 个体系。规范、组织、制度、人员管理。
          \item 6 大制度。侦查、检察、审判、监狱、律师、公证。
          \item 6 小制度。调解、陪审、死刑复核、再审、司法解释、案例指导。
          \item 司法解释审委会，指导案例可专委。
        \end{enumerate}
  \item 司法公正。最重要，最核心。\zd{针对的是某个案件}。
        \begin{enumerate}
          \item 实体。3 大结果，事实认定结果、法律适用结果、最终案件处理结果。
          \item 程序。过程、步骤、方法。
          \item 内涵。\zd{中参公平合廉正}。
                \begin{enumerate}
                  \item 司法人员的\zd{中}立性、\blh 程序的\zd{参}与性，各方的诉讼权利都能保障。
                  \item \blh 活动的\zd{公}开性。
                        \begin{enumerate}
                          \item 裁判文书上网，\zd{国}家\zd{秘}密、\zd{隐私}（离婚诉讼）、\zd{未成年、调解结案不公开}。
                          \item 检务公开。一般案件相关人员；重大社会人员。
                        \end{enumerate}
                  \item 当事人地位的\zd{平}等性、\blh 活动的\zd{合}法性。平等保护双方诉讼权利。
                  \item \blh 人员\zd{廉}洁性。
                        \begin{enumerate}
                          \item 开除吊证，终审禁止。
                          \item 禁止不当接触。3 工。\zd{工作需要}，\blh 时间，\blh 地点。报批；3 日内报备。
                        \end{enumerate}
                  \item \blh 结果的\zd{正}确性。事实清楚，证据确实充分。
                \end{enumerate}
          \item 审判独立与检察独立。独立行使职权。避免领导干部的干预；强化对司法权威的维护；加强对司法人员的履职保障。
        \end{enumerate}
  \item 效率。时间、资源、整体成本。公正优先，兼顾效率。
  \item 改革。政治性内容，涉及到方方面面的问题。
  \item 法律职业道德。不是很重要。
        \begin{enumerate}
          \item 特征。政治性、法律性、行业性（特殊的职业伦理）、专业性（资质）。
          \item 特点。职业性、实践性、正式性（规范文件）、严格性。没有国家强制性。还是道德层面。
          \item 6 大基本准则。看个 10 遍。常识判断。政治、正确、保密、配合、尽职、清廉。
        \end{enumerate}
\end{enumerate}
